%% Commands for TeXCount
%TC:macro \cite [option:text,text]
%TC:macro \citep [option:text,text]
%TC:macro \citet [option:text,text]
%TC:envir table 0 1
%TC:envir table* 0 1
%TC:envir tabular [ignore] word 
%TC:envir displaymath 0 word
%TC:envir math 0 word
%TC:envir comment 0 0
%%
%% The first command in your LaTeX source must be the \documentclass
%% command.
%%
%% For submission and review of your manuscript please change the
%% command to \documentclass[manuscript, screen, review]{acmart}.
%%
%% When submitting camera ready or to TAPS, please change the command
%% to \documentclass[sigconf]{acmart} or whichever template is required
%% for your publication.
%%
%%
\documentclass[sigconf, screen]{acmart}
%%
%% \BibTeX command to typeset BibTeX logo in the docs
\AtBeginDocument{%
  \providecommand\BibTeX{{%
    Bib\TeX}}}

%%
%% Additional packages required for this document
\usepackage{lipsum} % For generating dummy text
\usepackage{stfloats}
\usepackage{minted} 
\usepackage{units}
\usepackage{adjustbox}
\usepackage{hyperref}

%% Rights management information.  This information is sent to you
%% when you complete the rights form.  These commands have SAMPLE
%% values in them; it is your responsibility as an author to replace
%% the commands and values with those provided to you when you
%% complete the rights form.
\setcopyright{acmlicensed}
\copyrightyear{2025}
\acmYear{2025}
\acmDOI{XXXXXXX.XXXXXXX}
%% These commands are for a PROCEEDINGS abstract or paper.
\acmConference[RTKS '25]{Real-Time Kernels and Systems Exam}{Fill in Date}{Padova, Italy}
%%
%%  Uncomment \acmBooktitle if the title of the proceedings is different
%%  from ``Proceedings of ...''!
%%
%%\acmBooktitle{Woodstock '18: ACM Symposium on Neural Gaze Detection,
%%  June 03--05, 2018, Woodstock, NY}
\acmISBN{978-1-4503-XXXX-X/2018/06}


%%
%% Submission ID.
%% Use this when submitting an article to a sponsored event. You'll
%% receive a unique submission ID from the organizers
%% of the event, and this ID should be used as the parameter to this command.
%%\acmSubmissionID{123-A56-BU3}

%%
%% For managing citations, it is recommended to use bibliography
%% files in BibTeX format.
%%
%% You can then either use BibTeX with the ACM-Reference-Format style,
%% or BibLaTeX with the acmnumeric or acmauthoryear sytles, that include
%% support for advanced citation of software artefact from the
%% biblatex-software package, also separately available on CTAN.
%%
%% Look at the sample-*-biblatex.tex files for templates showcasing
%% the biblatex styles.
%%

%%
%% The majority of ACM publications use numbered citations and
%% references.  The command \citestyle{authoryear} switches to the
%% "author year" style.
%%
%% If you are preparing content for an event
%% sponsored by ACM SIGGRAPH, you must use the "author year" style of
%% citations and references.
%% Uncommenting
%% the next command will enable that style.
%%\citestyle{acmauthoryear}


%%
%% end of the preamble, start of the body of the document source.
\begin{document}

%%
%% The "title" command has an optional parameter,
%% allowing the author to define a "short title" to be used in page headers.
\title{Integration Report}

%%
%% The "author" command and its associated commands are used to define
%% the authors and their affiliations.
%% Of note is the shared affiliation of the first two authors, and the
%% "authornote" and "authornotemark" commands
%% used to denote shared contribution to the research.
\author{Gianluca Bresolin}
\email{gianluca.bresolin@studenti.unipd.it}
\affiliation{%
  \institution{Università degli Studi di Padova}
  \city{Padova}
  \state{Veneto}
  \country{Italy}
}

\author{Tommaso Prandin}
\email{tommaso.prandin@studenti.unipd.it}
\affiliation{%
  \institution{Università degli Studi di Padova}
  \city{Padova}
  \state{Veneto}
  \country{Italy}
}

%%
%% By default, the full list of authors will be used in the page
%% headers. Often, this list is too long, and will overlap
%% other information printed in the page headers. This command allows
%% the author to define a more concise list
%% of authors' names for this purpose.
%\renewcommand{\shortauthors}{Trovato et al.}

%%
%% The abstract is a short summary of the work to be presented in the
%% article.
\begin{abstract}
  This integration report includes some additional details regarding some
  aspects of the project that were not covered in the \textit{Techincal Report}.
\end{abstract}

%%
%% The code below is generated by the tool at http://dl.acm.org/ccs.cfm.
%% Please copy and paste the code instead of the example below.
%%
\begin{CCSXML}
<ccs2012>
   <concept>
       <concept_id>10011007.10010940.10010971.10011679</concept_id>
       <concept_desc>Software and its engineering~Real-time systems software</concept_desc>
       <concept_significance>500</concept_significance>
       </concept>
   <concept>
       <concept_id>10011007.10010940.10010971.10010564</concept_id>
       <concept_desc>Software and its engineering~Embedded software</concept_desc>
       <concept_significance>500</concept_significance>
       </concept>
   <concept>
       <concept_id>10011007.10010940.10010971.10010972</concept_id>
       <concept_desc>Software and its engineering~Software architectures</concept_desc>
       <concept_significance>300</concept_significance>
       </concept>
 </ccs2012>
\end{CCSXML}

\ccsdesc[500]{Software and its engineering~Real-time systems software}
\ccsdesc[500]{Software and its engineering~Embedded software}
\ccsdesc[300]{Software and its engineering~Software architectures}

%%
%% Keywords. The author(s) should pick words that accurately describe
%% the work being presented. Separate the keywords with commas.
\keywords{Real-Time Kernels, Real-Time Systems, Embedded Software, Software Architectures}
%% A "teaser" image appears between the author and affiliation
%% information and the body of the document, and typically spans the
%% page.
% \begin{teaserfigure}
%   \includegraphics[width=\textwidth]{example-image-a}
%   \caption{Caption}
%   \Description{Enjoying the baseball game from the third-base
%   seats. Ichiro Suzuki preparing to bat.}
%   \label{fig:teaser}
% \end{teaserfigure}

% \received{20 September 2025}
% \received[revised]{}
% \received[accepted]{}

%%
%% This command processes the author and affiliation and title
%% information and builds the first part of the formatted document.
\maketitle

\section{Why we used a one-to-one mapping between tasks and watchdogs in the Rust System?}
Since \textit{RTIC framework} does not provide any built-in mechanism to 
handle task deadlines, we decided to use a one-to-one mapping between
tasks and watchdogs. This way, each watchdog is responsible for monitoring
the deadline of a single task, making the implementation simpler and more
straightforward. \\
Besides the simplicity of implementation, this approach was chosen to avoid the 
runtime cost of managing a pool of deadlines and 
\texttt{deadline\_protected\_object}s at each deadline check.  
In fact, the alternative scenario of having a single watchdog
monitoring multiple tasks would require to identify over a pool of deadlines, at
each deadline check, the earliest next deadline in order to delay its execution
accordingly. This operation is needed also every task activation, since it could
introduce a new earliest deadline. \\
Even if we could optimize this operation by using an appropriate data structure,
such as a min-heap, the one-to-one mapping approach, instead, allows to manage
each watchdog independently, without the need of any search, insert or
delete operation over a pool of objects. This results in a constant-time
complexity for each deadline check and task activation, thus reducing
the runtime overhead, improving the system's responsiveness and reducing
watchdogs' intereference over system tasks. \\
Moreover, since the \textit{RTIC framework} requires to declare for each task
locals and shared resources at compile time, by using a one-to-one mapping we
can leverage the single responsability principle, enhancing code readability and
maintanability against a more complex single watchdog implementation. \\
However, this approach introduces a higher activation cost, since each watchdog
consists in a task, and as such it has its own activation overhead. This is a
price that we considered more accettable than the previous described runtime
cost, since the activation cost is paid only once and does not affect the system 
responsiveness. \\
Activation cost is not the only negative aspect of this approach: by having a
one-to-one mapping, we also introduce a memory overhead, since each watchdog
task requires its own stack and context. This could be a problem in systems with
a large number of tasks, where the memory overhead could become significant. In
our case, given the limited number of tasks in the system, we considered this
overhead accettable, but we acknowledge that in systems with a larger number of
tasks, this could become a trade-off to consider carefully against the runtime
cost of a single watchdog implementation.


\bibliographystyle{ACM-Reference-Format}
\bibliography{bibliography}


%%
%% If your work has an appendix, this is the place to put it.
% \appendix

\end{document}
\endinput
%%
